
\section{Gorizontal otish}

\begin{enumerate}

\item Tosh gorizontal otilgan. Otishdan 2 s o‘tib uning tezlik vektori gorizont bilan $45^\circ$ burchak hosil qildi. Toshning boshlang‘ich tezligini toping. $g=10\ \text{m/s}^2$.

\begin{enumerate}[label=\Alph*)]
\item 10
\item 20
\item 30
\item 40
\end{enumerate}

\item Jism gorizontal otilgan. Otishdan 2 s o‘tib tezlik va tezlanish vektorlarining yo‘nalishlari orasidagi burchak $60^\circ$ ga teng bo‘ldi. Jismning shu paytdagi to‘liq tezligini aniqlang. $g=10\ \text{m/s}^2$.

\begin{enumerate}[label=\Alph*)]
\item 20
\item 30
\item 40
\item 50
\end{enumerate}

\item Gorizontal yo‘nalishda otilgan jismning uchish uzoqligi u tashlangan balandlikning yarmiga teng. Jism yerga tushishida uning tezligi gorizont bilan hosil qilgan burchakning tangensi qanchaga teng?

\begin{enumerate}[label=\Alph*)]
\item 2
\item 4
\item 6
\item 8
\end{enumerate}

\item Tosh 5.5 m balandlikdan gorizontal yo‘nalishda shunday otildiki, u yer sirtiga $45^\circ$ burchak ostida tushadi. Tosh gorizontal yo‘nalishda qancha masofa uchib o‘tadi?

\begin{enumerate}[label=\Alph*)]
\item 11
\item 15
\item 22
\item 5.5
\end{enumerate}

\item 500 m balandlikda 180 km/h tezlik bilan uchayotgan samolyotdan yuk tushib ketdi. Qanday balandlikda yukning tezligi gorizontga $60^\circ$ burchak ostida yo‘nalgan bo‘ladi? $g=10\ \text{m/s}^2$.

\begin{enumerate}[label=\Alph*)]
\item 100
\item 125
\item 250
\item 375
\end{enumerate}

\item Koptok 2 m/s tezlik bilan gorizontal otildi. Koptokning gorizontal sirtga ikkita ketma-ket urilishlari orasidagi masofa 4 m ga teng. Koptok qanday balandlikdan otilgan? $g=10\ \text{m/s}^2$.

\begin{enumerate}[label=\Alph*)]
\item 2
\item 5
\item 10
\item 20
\end{enumerate}

\item Qiyalik burchagi gorizontga $30^\circ$ bo‘lgan tog‘dan 15 m/s boshlang‘ich tezlik bilan gorizontal yo‘nalishda koptok otiladi. Otish nuqtasidan qiya tekislik bo‘ylab qanday masofada koptok tushadi?

\begin{enumerate}[label=\Alph*)]
\item 15
\item 20
\item 30
\item 45
\end{enumerate}

\item Metall sharcha zinapoyaning yuqori zinasi chetiga 1.5 m/s tezlik bilan dumalab keldi. Har bir zinaning balandligi va kengligi 20 cm. Sharcha hisob bo‘yicha qaysi zinaga birinchi uriladi?

\begin{enumerate}[label=\Alph*)]
\item 2
\item 3
\item 4
\item 5
\end{enumerate}

\item Qo‘zg‘almas miltiqdan 120 m masofadagi vertikal nishonga gorizontal yo‘nalishda ikki marta o‘q uzilgan. Birinchi o‘qning tezligi 300 m/s, ikkinchisiniki 400 m/s. Nishondagi teshiklar orasidagi masofani (cm) aniqlang.

\begin{enumerate}[label=\Alph*)]
\item 25
\item 30
\item 35
\item 40
\end{enumerate}

\item Bir nuqtadan bir vaqtda ikki jism otiladi: biri 6 m/s tezlik bilan gorizontal, boshqasi 8 m/s tezlik bilan vertikal. Jismlar 2 s dan so‘ng bir-biridan qanday masofada bo‘ladi?

\begin{enumerate}[label=\Alph*)]
\item 10
\item 16
\item 18
\item 20
\end{enumerate}
\end{enumerate}


\begin{center}
\begin{tabular}{|c|c|c|c|c|}
\hline
1 & 2 & 3 & 4& 5 \\
\hline
B & C & B & A & B \\
\hline
6 & 7 & 8 & 9 & 10 \\
\hline
B & C & B & C & D \\
\hline
\end{tabular}
\end{center}





\section{Burchak ostida otish}

\begin{enumerate}

\item Koptok yer sirtidan gorizontga $30^\circ$ burchak ostida 20 m/s boshlang‘ich tezlik bilan otildi. U yerga urilguncha qancha vaqt o‘tadi?

\begin{enumerate}[label=\Alph*)]
\item 1
\item 2
\item 3
\item 4
\end{enumerate}

\item 40 m/s tezlik bilan gorizontga $60^\circ$ burchak ostida uchirilgan raketaning ko‘tarilish balandligini toping.

\begin{enumerate}[label=\Alph*)]
\item 20
\item 40
\item 60
\item 80
\end{enumerate}

\item Quroldan gorizontga burchak ostida uchib chiqqan snaryad 20 s uchdi. U qanday maksimal balandlikka erishgan?

\begin{enumerate}[label=\Alph*)]
\item 250
\item 400
\item 500
\item 600
\end{enumerate}

\item Gorizontga burchak ostida otilgan tosh 20 m maksimal balandikka erishdi. Toshning uchish vaqtini toping.

\begin{enumerate}[label=\Alph*)]
\item 2
\item 4
\item 6
\item 8
\end{enumerate}

\item 40 m/s tezlik bilan gorizontga $15^\circ$ burchak ostida uchirilgan raketaning uchish uzoqligini toping.

\begin{enumerate}[label=\Alph*)]
\item 40
\item 60
\item 80
\item 100
\end{enumerate}

\item Jism yer sirtidan gorizontga $30^\circ$ burchak ostida otilgan. To‘liq uchish vaqti 2 s ga teng bo‘ldi. Jismning boshlang‘ich tezligini toping.

\begin{enumerate}[label=\Alph*)]
\item 10
\item 15
\item 20
\item 25
\end{enumerate}

\item Ko‘tarilishning eng yuqori nuqtasida jismning tezligi boshlang‘ich tezlikdan ikki marta kichik bo‘lishi uchun uni gorizontga qanday burchak ostida otish kerak?

\begin{enumerate}[label=\Alph*)]
\item $30^\circ$
\item $45^\circ$
\item $60^\circ$
\item $75^\circ$
\end{enumerate}

\item Gorizontga burchak ostida otilgan tosh yerga 15 m/s tezlikda tushdi. Agar harakat vaqtida toshning eng katta tezligi eng kichik tezligidan uch marta katta bo‘lgani ma’lum bo‘lsa, uning maksimal ko‘tarilish balandligi qanday?

\begin{enumerate}[label=\Alph*)]
\item 5
\item 10
\item 15
\item 20
\end{enumerate}

\item 2 km balandlikda 540 km/h tezlik bilan uchayotgan samolyot aniq zenit quroli ustida bo‘lgan paytda quroldan o‘q uziladi. Snaryad qanday minimal tezlik bilan uchib chiqqanda nishonni shikastlashi mumkin?

\begin{enumerate}[label=\Alph*)]
\item 150
\item 200
\item 250
\item 300
\end{enumerate}

\item Gorizontga $45^\circ$ burchak ostida otilgan tosh harakat boshidan 0.8 s o‘tgach 12 m/s vertikal tashkil etuvchi tezlikka ega bo‘lgan holda ko‘tarilishni davom ettirdi. Toshning otilish va tushish nuqtalari orasidagi masofa qanday?

\begin{enumerate}[label=\Alph*)]
\item 40
\item 60
\item 80
\item 100
\end{enumerate}

\item Jism 30 m/s boshlang‘ich tezlik bilan gorizontga $45^\circ$ burchak ostida otildi. Jismning tezligi gorizontga $30^\circ$ burchak ostida yo‘nalgan paytda u qanday balandlikda bo‘ladi?

\begin{enumerate}[label=\Alph*)]
\item 10
\item 15
\item 20
\item 25
\end{enumerate}

\item Gorizontga $45^\circ$ burchak ostida otilgan disk 15 m maksimal balandlikka erishdi. Diskning uchish uzoqligi qanday?

\begin{enumerate}[label=\Alph*)]
\item 30
\item 45
\item 60
\item 75
\end{enumerate}

\item Tosh gorizontga shunday burchak ostida otildiki, bu burchakning sinusi 0.8 ga teng. Uchish uzoqligining maksimal ko‘tarilish balandligiga nisbatini toping.

\begin{enumerate}[label=\Alph*)]
\item 2
\item 3
\item 4
\item 5
\end{enumerate}

\item Toshni yer sirtidan birinchi gal gorizontga $60^\circ$ burchak ostida, ikkinchi gal $30^\circ$ burchak ostida otishdi. Agar birinchi holda tosh otilish joyidan ikki marta uzoqroqqa borib tushsa, birinchi holdagi ko‘tarilish balandligi ikkinchi holdagidan necha marta katta bo‘ladi?

\begin{enumerate}[label=\Alph*)]
\item 2
\item 4
\item 6
\item 8
\end{enumerate}

\item Yer sirtining bitta nuqtasidan ikkita tosh otib yuborildi. Birinchisi yerga $L$ masofada tushdi, ikkinchisi $3L$ masofada. Agar ikkinchi tosh gorizontga $30^\circ$ burchak ostida otilgan va ularning ko‘tarilish balandliklari birday bo‘lsa, birinchi tosh gorizontga qanday burchak ostida otilgan?

\begin{enumerate}[label=\Alph*)]
\item $30^\circ$
\item $45^\circ$
\item $60^\circ$
\item $75^\circ$
\end{enumerate}

\item Jism gorizontga burchak ostida otildi. Butun harakat vaqtining qanday qismida jism maksimal ko‘tarilish balandligining $\frac{3}{4}$ qismidan balandda bo‘ladi?

\begin{enumerate}[label=\Alph*)]
\item 25\%
\item 50\%
\item 75\%
\item 100\%
\end{enumerate}

\item Mina vertikal yuqoriga $3g$ tezlanish bilan uchadigan raketani shikastlashi uchun u raketa starti vaqtida minometdan qanday tezlikda uchib chiqishi kerak?

\begin{enumerate}[label=\Alph*)]
\item 50
\item 75
\item 100
\item 125
\end{enumerate}

\item Futbol to‘pi 10 m/s tezlik bilan gorizontga $15^\circ$ burchak ostida tepiladi. Tepish nuqtasidan 3 m masofada vertikal devor bo‘lib, to‘p unga elastik uriladi. To‘pning tushish masofasini toping.

\begin{enumerate}[label=\Alph*)]
\item 1
\item 2
\item 3
\item 4
\end{enumerate}

\item Shlangdan gorizontga $15^\circ$ burchak ostida suv oqimi chiqadi. Oqim shlangdan 20 m masofada borib tushadi. Shlangdan 1 minutda qancha massali suv oqib chiqadi?

\begin{enumerate}[label=\Alph*)]
\item 60
\item 90
\item 120
\item 150
\end{enumerate}

\item Shlangdan 10 m/s tezlik bilan gorizontga $30^\circ$ burchak ostida suv oqimi chiqadi. Havoda bo‘ladigan suvning massasini toping.

\begin{enumerate}[label=\Alph*)]
\item 1
\item 2
\item 3
\item 4
\end{enumerate}

\item 15 m balandlikda joylashgan nuqtadan tosh 20 m/s tezlik bilan gorizontga $30^\circ$ burchak ostida otiladi. Qancha vaqtdan so‘ng u yerga tushadi?

\begin{enumerate}[label=\Alph*)]
\item 2
\item 3
\item 4
\item 5
\end{enumerate}

\item Yerdan 30 m balandlikda joylashgan nuqtadan tosh 20 m/s tezlik bilan gorizontga $45^\circ$ burchak ostida otiladi. U qanday masofada tushadi?

\begin{enumerate}[label=\Alph*)]
\item 40
\item 50
\item 60
\item 70
\end{enumerate}

\item 7.5 m balandlikda joylashgan derazadan tosh gorizontga $45^\circ$ burchak ostida otiladi. U uy devoridan 15 m masofada tushadi. Tosh qanday tezlik bilan otilgan?

\begin{enumerate}[label=\Alph*)]
\item 5
\item 10
\item 15
\item 20
\end{enumerate}

\item Tog‘ yonbag‘ridagi nuqtadan 21 m/s boshlang‘ich tezlik bilan gorizontga $60^\circ$ burchak ostida jism otiladi. Jism qanday masofada tushadi?

\begin{enumerate}[label=\Alph*)]
\item 15
\item 20
\item 30
\item 45
\end{enumerate}

\item Koptok qiyalik burchagi gorizontga $30^\circ$ bo‘lgan tog‘ ustida 6 m/s boshlang‘ich tezlik bilan tog‘ yonbag‘riga tik ravishda otib yuboriladi. Koptok qanday masofada tushadi?

\begin{enumerate}[label=\Alph*)]
\item 240
\item 360
\item 480
\item 600
\end{enumerate}

\item 1.5 m balandlikdan qiya tekislik ustiga vertikal yo‘nalishda sharcha tushadi va absolyut elastik ravishda qaytadi. U yana shu tekislikka tushish nuqtasidan qanday masofada uriladi?

\begin{enumerate}[label=\Alph*)]
\item 2
\item 4
\item 6
\item 8
\end{enumerate}

\item Ikkita tosh bir vaqtda otiladi: biri shimoliy yo‘nalishda 24 m/s tezlik bilan, boshqasi janubiy yo‘nalishda 32 m/s tezlik bilan. 1.5 s dan so‘ng ular orasidagi masofani toping.

\begin{enumerate}[label=\Alph*)]
\item 40
\item 50
\item 60
\item 70
\end{enumerate}

\item Ikkita tosh bir-biridan 42 m masofada joylashgan. Harakat davomida toshlar orasidagi eng kichik masofa qanchaga teng?

\begin{enumerate}[label=\Alph*)]
\item 6
\item 8
\item 10
\item 12
\end{enumerate}

\end{enumerate}



\begin{center}
\begin{tabular}{|c|c||c|c|}
\hline
1 & B & 15 & C \\
\hline
2 & C & 16 & B \\
\hline
3 & C & 17 & C \\
\hline
4 & B & 18 & A \\
\hline
5 & C & 19 & C \\
\hline
6 & C & 20 & B \\
\hline
7 & C & 21 & B \\
\hline
8 & B & 22 & C \\
\hline
9 & C & 23 & B \\
\hline
10 & C & 24 & C \\
\hline
11 & B & 25 & C \\
\hline
12 & C & 26 & C \\
\hline
13 & B & 27 & C \\
\hline
14 & C & 28 & B \\
\hline
\end{tabular}
\end{center}









